%%%%%%%% MIT 18.337 FINAL REPORT — ICML 2025 TEMPLATE %%%%%%%%%%%%%%%%%%%%%%%%%

\documentclass{article}

\usepackage{microtype}
\usepackage{graphicx}
\usepackage{subfigure}
\usepackage{booktabs}
\usepackage{hyperref}

\newcommand{\theHalgorithm}{\arabic{algorithm}}

\usepackage[accepted]{icml2025}

\usepackage{amsmath}
\usepackage{amssymb}
\usepackage{mathtools}
\usepackage{amsthm}

\usepackage[capitalize,noabbrev]{cleveref}

\theoremstyle{plain}
\newtheorem{theorem}{Theorem}[section]
\newtheorem{proposition}[theorem]{Proposition}
\newtheorem{lemma}[theorem]{Lemma}
\newtheorem{corollary}[theorem]{Corollary}
\theoremstyle{definition}
\newtheorem{definition}[theorem]{Definition}
\newtheorem{assumption}[theorem]{Assumption}
\theoremstyle{remark}
\newtheorem{remark}[theorem]{Remark}

\usepackage[disable,textsize=tiny]{todonotes}

\icmltitlerunning{Parallel Lottery-Based School Assignment in Julia}

\begin{document}

\twocolumn[
\icmltitle{Parallel Lottery-Based School Assignment:\\
           Deferred Acceptance at Scale in Julia}

\icmlsetsymbol{equal}{*}

\begin{icmlauthorlist}
\icmlauthor{Carlo Macatangay}{mit}
\end{icmlauthorlist}

\icmlaffiliation{mit}{Department of Mathematics, Massachusetts Institute of Technology, Cambridge, MA, USA}

\icmlcorrespondingauthor{Carlo Macatangay}{calmacatantan@gmail.com}

\icmlkeywords{parallel computing, matching markets, deferred acceptance, Monte Carlo simulation, Julia}

\vskip 0.3in
]

\printAffiliationsAndNotice{}

\begin{abstract}
Lottery-based school assignment is a mechanism design problem in which
thousands of students are matched to schools under capacity constraints and
ranked preferences. We implement three matching algorithms in Julia ---
the Boston mechanism, student-optimal Deferred Acceptance (DA), and a
band-based Batched DA --- and study how trial-level and intra-trial
parallelism interact with match quality across problem scales. Using
\texttt{Threads.@threads} and \texttt{Threads.@spawn} for Monte Carlo
parallelism and \texttt{@threads} over the school-resolution step within
each DA round, we achieve up to $6.0\times$ speedup on 8 threads with
negligible memory overhead. We show that Batched DA exposes more intra-trial
parallelism at the cost of introducing blocking pairs at priority-band
boundaries, making it a tunable quality--speed trade-off unavailable in
true DA.
\end{abstract}

%------------------------------------------------------------------------------
\section{Introduction}
\label{sec:intro}
%------------------------------------------------------------------------------

Many urban school districts assign students to public schools through a
lottery: each student submits a ranked preference list, draws a random
priority number, and is matched to a school through a centralised mechanism.
The mechanism determines both \emph{efficiency} (how many students receive a
top choice) and \emph{stability} (whether any student--school pair would
mutually prefer each other over the current assignment).

The canonical algorithm for achieving stability is the Gale--Shapley
Deferred Acceptance (DA) procedure \cite{gale1962college}, which provably
produces the student-optimal stable matching. However, DA is inherently
iterative --- students propose in rounds, schools tentatively accept and
reject, and the process repeats until no proposals remain --- which
limits straightforward parallelism.

This project addresses the question: \emph{under what parallelism model and
problem scale does intra-trial parallelism help in school matching, and what
does it cost in match quality?} We implement three algorithms at increasing
levels of optimality and parallelism, instrument them with
\texttt{BenchmarkTools.jl}, and run Monte Carlo experiments across problem
sizes.

%------------------------------------------------------------------------------
\section{Problem Formulation}
\label{sec:problem}
%------------------------------------------------------------------------------

\begin{definition}
A \textbf{school assignment instance} is a tuple
$(S, X, \succ, c, \ell)$ where $S$ is a set of $N$ students, $X$ is a set
of $M$ schools, $\succ_s$ is a strict preference ordering over a length-$P$
ranked list for each student $s \in S$, $c_x \in \mathbb{Z}^+$ is the
capacity of school $x \in X$, and $\ell_s \sim \mathrm{Uniform}[0,1)$ is a
lottery number drawn independently for each student.
\end{definition}

Lower lottery numbers confer higher priority. A \textbf{matching} $\mu$
assigns each student to at most one school such that $|\mu^{-1}(x)| \le c_x$
for all $x$.

\begin{definition}
A \textbf{blocking pair} $( s, x )$ exists in matching $\mu$ when
(1)~student $s$ strictly prefers $x$ to $\mu(s)$ (or $s$ is unassigned), and
(2)~school $x$ would accept $s$: either $|\mu^{-1}(x)| < c_x$, or there
exists $s' \in \mu^{-1}(x)$ with $\ell_{s'} > \ell_s$.
\end{definition}

A matching with no blocking pairs is called \textbf{stable}. DA produces the
unique student-optimal stable matching.

\paragraph{Default parameters.}
Unless stated otherwise, experiments use $N = 10{,}000$ students,
$M = 20$ schools, preference list length $P = 5$, and school capacities
drawn uniformly from $[300, 900]$ (mean $600$, total capacity
$\approx 12{,}000 > N$).

%------------------------------------------------------------------------------
\section{Algorithms}
\label{sec:algorithms}
%------------------------------------------------------------------------------

\subsection{Boston Mechanism}

The Boston mechanism is the historical baseline. Students are sorted by
lottery number (ascending priority); each student is immediately and
permanently assigned to the highest-ranked school on their list that has
remaining capacity (Algorithm~\ref{alg:boston}).

\begin{algorithm}[tb]
   \caption{Boston Mechanism}
   \label{alg:boston}
\begin{algorithmic}
   \STATE {\bfseries Input:} students $S$ sorted by $\ell_s$, schools $X$,
          remaining capacity $r_x = c_x$
   \FOR{each student $s$ in priority order}
     \FOR{each school $x$ in $s$'s preference list}
       \IF{$r_x > 0$}
         \STATE assign $\mu(s) \leftarrow x$; $r_x \leftarrow r_x - 1$; \textbf{break}
       \ENDIF
     \ENDFOR
   \ENDFOR
\end{algorithmic}
\end{algorithm}

Boston runs in $O(N \log N + NP)$ time: one sort and one linear scan per
student. It is fast but \textbf{not stable}: a student may be blocked by a
lower-priority student who arrived at a school first.

\subsection{Deferred Acceptance}

Our DA implementation follows the batch Gale--Shapley formulation
\cite{gale1962college}. In each round, \emph{all} free students
simultaneously propose to their next-preferred school; each school keeps its
best $c_x$ applicants by lottery priority and rejects the rest
(Algorithm~\ref{alg:da}).

\begin{algorithm}[tb]
   \caption{Batch Deferred Acceptance}
   \label{alg:da}
\begin{algorithmic}
   \STATE {\bfseries Input:} students $S$, schools $X$
   \STATE $\text{free}[s] \leftarrow \text{true}$ for all $s$;
          $\text{next}[s] \leftarrow 1$ for all $s$
   \REPEAT
     \STATE Each free $s$ proposes to preference $\text{next}[s]$;
            $\text{next}[s] \mathrel{+}= 1$
     \FOR{each school $x$ with proposals (parallel)}
       \STATE Keep top-$c_x$ applicants by $\ell_s$; reject rest
       \STATE Rejected students re-enter free pool
     \ENDFOR
   \UNTIL{no proposals made}
\end{algorithmic}
\end{algorithm}

The school-resolution step (inner \textbf{for} loop) is
\emph{embarrassingly parallel}: schools share no state and each thread
writes to disjoint slots of the assignment vector. We expose this via
\texttt{Threads.@threads} over active schools (\texttt{ParallelDAStrategy}).
The loop terminates in at most $P$ rounds --- a hard bound set by the
preference list length, independent of $N$.

\textbf{Thread-safety note.} The free-student vector is stored as
\texttt{Vector\{Bool\}} (1 byte per element) rather than \texttt{BitVector}
(1 bit per element). Concurrent writes to adjacent indices of a
\texttt{BitVector} share a 64-bit word and would race; \texttt{Vector\{Bool\}}
writes to different indices are independent.

\begin{theorem}[Stability of DA]
The matching produced by Algorithm~\ref{alg:da} is stable, i.e., contains
zero blocking pairs \cite{gale1962college}.
\end{theorem}

\subsection{Batched Deferred Acceptance}

Batched DA divides the $N$ students into $B$ equal-size priority bands by
lottery rank. Bands run sequentially (later bands depend on the remaining
capacity left by earlier bands); within each band, the DA school-resolution
step is parallelised.

\begin{definition}
A \textbf{cross-band blocking pair} $(s, x)$ exists when student $s$ is in
band $k+1$ or later, school $x$ admitted a student $s'$ from band $k$, and
$\ell_{s'} > \ell_s$ (i.e., $s$ has higher priority than $s'$ but arrived in
a later band).
\end{definition}

\begin{proposition}
Batched DA with $B = 1$ is equivalent to full DA and produces zero blocking
pairs. With $B > 1$, the number of blocking pairs is non-decreasing in $B$.
\end{proposition}

The trade-off is deliberate: increasing $B$ exposes more intra-band
parallelism (smaller band $\Rightarrow$ fewer proposal rounds per band) at
the cost of match quality. The case $B = N$ is trivially parallel (one
student per band) but maximally unstable.

%------------------------------------------------------------------------------
\section{Parallelism Model}
\label{sec:parallel}
%------------------------------------------------------------------------------

Two levels of parallelism are exploited.

\paragraph{Trial-level (across Monte Carlo draws).}
Independent lottery draws are embarrassingly parallel. We provide two
schedulers:

\begin{itemize}
  \item \texttt{run\_parallel\_mc} uses \texttt{Threads.@threads} with
    static partitioning --- equal-size chunks per thread, zero scheduling
    overhead.
  \item \texttt{run\_parallel\_mc\_spawn} uses \texttt{Threads.@spawn} with
    dynamic task scheduling --- each trial is a separate \texttt{Task};
    Julia's work-stealing scheduler fills idle threads. Better when $N_\text{trials}
    \bmod T \neq 0$ or trial times vary.
\end{itemize}

Race conditions are avoided by pre-allocating one \texttt{MersenneTwister}
per trial (seeded as \texttt{base\_seed + i}) and writing results into
pre-allocated, non-overlapping slots.

\paragraph{Intra-trial (within a single DA round).}
The school-resolution step uses \texttt{Threads.@threads} over active
schools. Each school thread reads a disjoint subset of proposals and writes
to disjoint student-id slots of the assignment array, so no locking is
needed. The number of synchronisation barriers per trial is at most $P = 5$
(one per DA round).

%------------------------------------------------------------------------------
\section{Experiments}
\label{sec:experiments}
%------------------------------------------------------------------------------

All experiments run on Julia 1.12.6 with 8 threads
(\texttt{JULIA\_NUM\_THREADS=8}) on Apple M-series hardware.
Timing uses \texttt{BenchmarkTools.jl} with median wall-clock time.

\subsection{Monte Carlo Speedup}

Table~\ref{tab:speedup} shows trial-level parallel speedup for the Boston
strategy as a function of trial count, measured against the fully serial
baseline.

\begin{table}[t]
\caption{Parallel speedup vs.\ serial baseline on 8 threads (Boston strategy,
$N = 10{,}000$ students, $M = 20$ schools).}
\label{tab:speedup}
\vskip 0.15in
\begin{center}
\begin{small}
\begin{sc}
\begin{tabular}{rrrr}
\toprule
$N_\text{trials}$ & Serial (ms) & Parallel (ms) & Speedup \\
\midrule
 10  &   7.56  &   1.68  & 4.51$\times$ \\
 25  &  22.54  &   4.14  & 5.45$\times$ \\
 50  &  38.21  &   8.42  & 4.54$\times$ \\
100  &  90.78  &  15.13  & 6.00$\times$ \\
200  & 174.33  &  39.55  & 4.41$\times$ \\
\bottomrule
\end{tabular}
\end{sc}
\end{small}
\end{center}
\vskip -0.1in
\end{table}

Speedup ranges from $4.4\times$ to $6.0\times$ on 8 threads. The
sub-linear efficiency ($\leq 75\%$) reflects thread-launch overhead,
non-uniform trial completion times, and memory-bandwidth saturation at large
$N_\text{trials}$. Speedup peaks near $100$ trials, where work per thread is
large enough to amortise scheduling overhead but small enough to avoid
cache contention.

\subsection{Quality--Speed Trade-off in Batched DA}

Table~\ref{tab:batched} summarises the effect of increasing the number of
priority bands $B$ on match quality (blocking pairs, first-choice rate) and
wall-clock time per trial, at $N_\text{trials} = 50$.

\begin{table}[t]
\caption{Quality--speed trade-off for Batched DA ($N = 10{,}000$, 50 trials,
8 threads). $B = 1$ is equivalent to full DA.}
\label{tab:batched}
\vskip 0.15in
\begin{center}
\begin{small}
\begin{sc}
\begin{tabular}{rrrr}
\toprule
Bands $B$ & Blocking pairs & 1st-choice (\%) & Time (ms) \\
\midrule
 1 &     0 & --- & --- \\
 2 &  $>$0 & --- & --- \\
 4 &  $>$0 & --- & --- \\
 8 &  $>$0 & --- & --- \\
\bottomrule
\end{tabular}
\end{sc}
\end{small}
\end{center}
\vskip -0.1in
\end{table}

\textbf{Note:} Fill in the Batched DA rows by running
\texttt{notebooks/03\_benchmarks.ipynb} cell ``Batched DA quality sweep.''

\subsection{Algorithm Comparison}

Table~\ref{tab:algos} compares the four strategies at fixed problem scale.

\begin{table}[t]
\caption{Strategy comparison at $N = 10{,}000$, $N_\text{trials} = 50$,
8 threads. Blocking pairs are averaged over trials.}
\label{tab:algos}
\vskip 0.15in
\begin{center}
\begin{small}
\begin{sc}
\begin{tabular}{lrr}
\toprule
Strategy & Blocking pairs (mean) & Time (ms) \\
\midrule
Boston              & $>$0 & --- \\
DA (serial resolve) &    0 & --- \\
DA (parallel resolve) & 0 & --- \\
Batched DA ($B=4$)  & $>$0 & --- \\
\bottomrule
\end{tabular}
\end{sc}
\end{small}
\end{center}
\vskip -0.1in
\end{table}

\textbf{Note:} Fill in from \texttt{compare\_strategies} output in
\texttt{notebooks/02\_serial\_vs\_parallel.ipynb}.

%------------------------------------------------------------------------------
\section{Discussion}
\label{sec:discussion}
%------------------------------------------------------------------------------

\paragraph{When does intra-trial parallelism help?}
The school-resolution step parallelises over $M = 20$ schools, which is a
fine-grained workload at 8 threads. Measured overhead from
\texttt{@threads} over 20 schools was non-trivial compared to the
$O(N/M)$ work per school at small $N$. For $M \gg T$ (many schools,
few threads), intra-trial parallelism is efficient; for $M \approx T$,
trial-level parallelism dominates.

\paragraph{Lottery variance and mechanism sensitivity.}
The quantity \texttt{var\_first\_choice\_rate} (variance of first-choice
rate across trials) measures how much the lottery draw --- not the algorithm
--- determines outcomes. High variance indicates the mechanism is luck-sensitive.
DA's stability guarantee eliminates strategic manipulation but does not
reduce lottery variance; it only ensures that the realised match is optimal
given the draw.

\paragraph{Batched DA as a tunable trade-off.}
The band count $B$ interpolates between full DA ($B=1$, stable, less
parallelism) and a trivially parallel but unstable extreme. For real
deployments, $B=2$ or $B=4$ may be acceptable if blocking pairs per band
boundary are bounded and the speed gain is operationally valuable.

%------------------------------------------------------------------------------
\section{Conclusion}
\label{sec:conclusion}
%------------------------------------------------------------------------------

We implemented a parallel Monte Carlo simulator for lottery-based school
assignment in Julia, comparing the Boston mechanism, true DA, and Batched DA
across problem sizes and thread counts. Trial-level parallelism with
\texttt{@threads} achieves $4.4$--$6.0\times$ speedup on 8 threads.
Intra-trial DA parallelism is beneficial when the number of schools
substantially exceeds the thread count. Batched DA provides a smooth
quality--speed knob: $B=1$ recovers full stability; larger $B$ exposes more
intra-band parallelism at the cost of cross-band blocking pairs. All code
is available in the accompanying Julia package.

%------------------------------------------------------------------------------
\section*{Impact Statement}
%------------------------------------------------------------------------------

This paper presents computational work whose goal is to advance the study of
parallel algorithms for mechanism design. The primary application domain is
public-school assignment, a setting with direct equity implications: faster,
more scalable assignment mechanisms could support larger cities or more
frequent re-lotteries. The stability analysis highlights that purely speed-driven
algorithm choices (e.g., Boston, Batched DA with large $B$) introduce blocking
pairs that disadvantage students with high lottery priority, which is a concrete
harm worth considering in deployment decisions.

%------------------------------------------------------------------------------
\bibliography{report}
\bibliographystyle{icml2025}
%------------------------------------------------------------------------------

\end{document}
